\section{Moving frames, cyclic covers, and residues}
\label{bl:cov:section}
\label{bl:sec:covers}

The first step in Proposition~\ref{bl:prop:complex-supported} is to
embed the entire lifting obstruction into logarithmic residue
cohomology. Moving to a frame bundle supplies compatible roots without
choosing roots separately on open sets. The complete cyclic covers
retain the character summand needed for an injective residue map,
including the poles along the extra coefficient-one boundary.

Work over $\C$, and put $A=K_V+C$ and $S=G_{\mathrm{red}}$.
Replace $(q,G)$ by
$(bq,bG)$, for a sufficiently divisible positive integer $b$, so that
semiampleness on the entire reduced scheme $S$ gives fixed identifications
\begin{equation}\label{bl:cov:initial-data}
 L=\OO_V(G)\simeq\OO_V(qA),\qquad
 L|_S\simeq f^*J,\qquad
 f:S\longrightarrow T=\PP^\ell,\quad J=\OO_T(1).
\end{equation}
Indeed, a globally generated power of the original restriction defines
$f$ by its complete linear system; surjectivity onto $T$ is unnecessary.
Further powers preserve this property.  We also clear all Cartier indices
and boundary denominators that occur below.  Let $s\in H^0(V,L)$ be the
section with divisor $G$, and write
\[
 G=\sum_i d_iS_i,\qquad S=\sum_iS_i,\qquad C=S+H+C_0.
\]
Here $d_i>0$, $H$ is the reduced sum of the remaining coefficient-one
components, and $C_0$ has coefficients in $[0,1)$ and shares no component
with $S+H$.  Choose $r$ divisible by $q$ and all $d_i$, and put $a=r/q$.
The singularity properties used in this section, and in the remainder of
the supported-pair argument, are log canonicity and klt singularities away
from $S+H$.  A dlt pair has both properties.  Apart from normality, the
Cartier hypotheses used are that $K_V$ and the indicated boundary
components are $\Q$-Cartier, together with the actual identifications
in~\eqref{bl:cov:initial-data}.

\subsection{The moving frame and the first cover}

Let $p:P=L^\times\to V$ and $B=J^\times\to T$ denote the bundles of
nonzero vectors.  The tautological frame $w$ trivializes $p^*L$ on $P$.
For $\lambda\in\Gm$, let $\lambda$ multiply the vectors in both bundles
by $\lambda^r$.  All weights below are pullback weights: a tensor $u$
has weight $m$ if $\lambda^*u=\lambda^m u$.  In particular $w$ has
weight $r$.  The fixed restriction isomorphism gives
\[
 S_P:=p^{-1}(S)=S\times_TB\longrightarrow B,
\]
an equivariant projective morphism.

There is an invariant canonical comparison
$\omega_P\simeq p^*\omega_V$, understood reflexively on the normal
spaces.  On a bundle chart with nonzero fiber coordinate $t$, it sends a
top form $\eta$ on $V$ to $p^*\eta\wedge dt/t$.  Replacing $t$ by
$ut$, with $u$ invertible on the base, adds $du/u$ to $dt/t$; its wedge
with $\eta$ vanishes.  Thus the comparison glues and is invariant under
scaling.  The same construction applies to divisible adjoint powers.
Moreover, the pullback pair $(P,p^*C)$ is lc and is klt off $S_P+p^*H$:
locally a log resolution is the product of a log resolution on $V$ with
$\Gm$, with the same discrepancy coefficients.

The ratio $u=p^*s/w$ is a regular function on $P$.  Form the whole finite
algebra $\OO_P[y]/(y^r-u)$ and let
\begin{equation}\label{bl:cov:first-cover}
 \pi:Z\longrightarrow P
\end{equation}
be its normalization. This is the cyclic-cover construction of
\cite[Section~3.5]{EsnaultViehweg1992}, applied to the trivial line
bundle; we retain its full-algebra convention. Throughout, normalization of a generically
reduced algebra means normalization in its full total ring of fractions.
Thus $Z$ is a disjoint union of normal varieties if the generic algebra
splits; no component is discarded.  The polynomial is separable over
every generic field, and normalization is finite since the schemes are
of finite type over $\C$.  The function $y$ remains regular, and the
action lifts by
\[
 \lambda^*y=\lambda^{-1}y.
\]
It commutes with the deck action $y\mapsto\zeta y$, $\zeta\in\mu_r$.

\begin{lemma}\label{bl:cov:reduced-cover}
The divisor $E=(y=0)$ on $Z$ is reduced and Cartier, with ideal
$I=(y)=\OO_Z(-E)$.  Put
$H_Z=\pi^*p^*H$ and $C_{0,Z}=\pi^*p^*C_0$.  With compatible canonical
divisors one has
\begin{equation}\label{bl:cov:crepant}
 K_Z+E+H_Z+C_{0,Z}=\pi^*(K_P+p^*C).
\end{equation}
The pair $(Z,E+H_Z+C_{0,Z})$ is lc and is klt off $E+H_Z$.
The map $E\to S_P$ induces an equivariant projective morphism
$g:E\to B$.
\end{lemma}

\begin{proof}
At the generic point of $p^{-1}(S_i)$, write $u=vx^{d_i}$ in the
corresponding discrete valuation ring, with $v$ a unit.  After an
unramified extension extracting a root of $v$, the normalized Kummer
algebra has ramification index $e_i=r/d_i$ on each branch, since $d_i$
divides $r$.  Consequently
\[
 \ord_F(y)=e_id_i/r=1
\]
at every prime $F$ lying over this divisor.  These are all the zero
divisors of $y$.  On every normal component of $Z$, the nonzero function
$y$ is a nonzerodivisor and its principal ideal is radical.  For the
latter assertion, if $v^N\in(y)$, the height-one valuation inequalities
give $\ord_F(v)\geq\ord_F(y)$ at every prime; normality then implies
$v/y\in\OO_Z$.  This also proves that $E$ is reduced as a scheme.

Off $S_P$ the first root algebra is finite \'{e}tale.  At a prime above
$S_P$ the canonical ramification coefficient is $e_i-1$, whereas the
pullback coefficient of $S_P$ is $e_i$.  This proves
\eqref{bl:cov:crepant} in codimension one, and hence as an identity of
$\Q$-Cartier divisors.  In particular $H_Z$ is reduced and the
coefficients of $C_{0,Z}$ are those of $C_0$.

We record why this calculation controls all divisorial valuations,
including exceptional ones.  For a finite dominant map of normal
varieties $X'\to X$ in characteristic zero and a comparison
$K_{X'}+\Delta'=\phi^*(K_X+\Delta)$, let $v'$ be a normalized divisorial
valuation upstairs.  Its restriction is $e v$ for a normalized
divisorial valuation $v$ downstairs: the residue extension is finite,
so its transcendence degree is still $\dim X-1$.  Choose a model of $X$
on which $v$ is a divisor and normalize it in the upstairs function
field.  At the resulting discrete valuation rings the different has
order $e-1$ (the extension is tame).  Comparing the canonical
coefficients on these models yields
\begin{equation}\label{bl:cov:finite-discrepancy}
 A_{X',\Delta'}(v')=e A_{X,\Delta}(v),
\end{equation}
where $A$ denotes log discrepancy.  The same argument applies to each
field factor of a disconnected cover.  Applied to
\eqref{bl:cov:crepant}, it gives lc singularities everywhere and klt
singularities over $P\setminus(S_P+p^*H)$.  Finally, the reduced scheme
$E$ maps into $S_P$ because its image is supported there.  The resulting
map is finite, so its composite with $S_P\to B$ is projective.
\end{proof}

\subsection{The root of the adjoint tensor}

The nowhere-vanishing frame $w$, the fixed isomorphism in
\eqref{bl:cov:initial-data}, and the canonical comparison produce a
meromorphic $q$-canonical tensor on $P$.  As an adjoint section it has
no zeros, so its divisor as a meromorphic tensor is $-qp^*C$.
Its differential pullback to $Z$ is a tensor $\tau$, nonzero on every
component, satisfying
\begin{equation}\label{bl:cov:tau}
 \Div(\tau)=-q\Delta_Z,\qquad
 \Delta_Z=E+H_Z+C_{0,Z},\qquad \wt(\tau)=r.
\end{equation}
The canonical ramification formula used in
\eqref{bl:cov:crepant} proves this equality, including its integral
coefficients.  Thus the construction depends on the given rational
linear equivalence, not merely on a numerical equivalence.

Let $\nu:\widetilde Z\to Z$ be the full normalized $q$-th root of
$\tau$.  Explicitly, for a meromorphic canonical frame $\theta$ on a
component with function field $K$, normalize in the finite reduced
$K$-algebra
\begin{equation}\label{bl:cov:second-algebra}
 K[c]/\bigl(c^q-\tau/\theta^{\otimes q}\bigr).
\end{equation}
Changing $\theta$ to $v\theta$ changes $c$ to $c/v$, so the algebras
and the tautological form $c\theta$ glue.  Pullbacks of forms in this
description are differential pullbacks.  The deck group $\mu_q$ acts
by $c\mapsto\zeta c$, and the tautological form has its character
$\chi(\zeta)=\zeta$.  Even if the algebra splits, its invariant
subalgebra is exactly $K$: the basis $1,c,\ldots,c^{q-1}$ has all the
distinct characters.  This observation will be needed for descent of
ratios of forms.

The $\Gm$-action also lifts algebraically.  If on a meromorphic chart
$\lambda^*\theta=j_\lambda\theta$, using differentials relative to the
parameter $\lambda$, prescribe
\begin{equation}\label{bl:cov:second-action}
 \lambda^*c=\lambda^a j_\lambda^{-1}c.
\end{equation}
This respects the equation because $aq=r$ and
$\lambda^*\tau=\lambda^r\tau$; the factors $j_\lambda$ satisfy the
cocycle identity.  These formulas define the action on the generic
algebra and on its integral closure.  More explicitly, normalization
commutes with the smooth base extension by $\Gm$, and the isomorphism
between the two pullback algebras over $\Gm\times Z$ therefore induces
an isomorphism of their normalizations.  This gives a regular action,
not just a separate lift of each scalar automorphism.  It commutes
with $\mu_q$, and the tautological form has weight $a$.

Take a projective $(\Gm\times\mu_q)$-equivariant log resolution of
$\widetilde Z$ and the inverse images of the boundary supports, and
write its composite with $\nu$ as
\[
 \rho:Y\longrightarrow Z.
\]
Functorial resolution and principalization in characteristic zero
give this resolution by invariant centers and projective blowups
\cite[Theorem~1.1]{BierstoneMilman2008},
\cite[Theorem~1.0.1]{Wlodarczyk2005}.
Indeed, functoriality for the smooth action and projection morphisms
$(\Gm\times\mu_q)\times\widetilde Z\to\widetilde Z$ lifts the action
and its group identities at every step.
We use the pulled-back tautological top form
$\sigma$ on $Y$; it has deck character $\chi$ and weight $a$.
We may choose $Y$ smooth and quasi-projective with a $\Gm$-linearized
ample line bundle.  To see the latter point, compactify $P$ as
$\PP_V(\OO_V\oplus L)$, with its scaling action.  A sufficiently
positive twist of the tautological bundle by an ample bundle from
$V$ is ample and linearized.  Restrict to $P$ and pull back through
the finite covers, which preserves ampleness and linearization.
For every invariant blowup center its ideal, and hence the
tautological relatively ample bundle on the blowup, is linearized.
Tensoring these with sufficiently high powers of the preceding ample
bundles at successive steps gives the assertion for $Y$.

Define effective divisors on $Y$ by
\begin{equation}\label{bl:cov:resolved-boundary}
 \widehat E=\rho^*E,\qquad
 D_h=(\Supp\rho^*H_Z)_{\mathrm{red}},\qquad
 R=\sum_{\substack{F\text{ prime in }\Supp\widehat E\\
                         F\not\subset D_h}} F.
\end{equation}
The union $R+D_h$ is reduced simple normal crossing.  Both $R$ and
$\widehat E$ map scheme-theoretically to $E$: the function $y$ vanishes
on them.  Write $q_R:R\to E$ and $q_{\widehat E}:\widehat E\to E$ for
these projective maps, and set $h=g\circ q_R:R\to B$.

The spaces and morphisms are summarized in Figure~\ref{fig:boundary-covers}.
The form $\sigma$ lives on $Y$; its residue will live on $R$ and will
be pushed through the top row to $B$.
\begin{figure}[ht]
\centering
\begin{tikzcd}[column sep=large,row sep=large]
 R \arrow[r,"q_R"] \arrow[d,hook] &
 E \arrow[r] \arrow[d,hook] &
 S_P \arrow[r] \arrow[d,hook] & B\\
 Y \arrow[r,"\rho"'] & Z \arrow[r,"\pi"'] & P &
\end{tikzcd}
\caption{The reduced boundary maps to $B$ through the top row, whose
composite is $h$. The lower row contains the first normalized root
cover $\pi$ and the second normalized cover followed by resolution,
combined in $\rho$. Only the boundary is equipped with a map to $B$.}
\label{fig:boundary-covers}
\end{figure}

\begin{lemma}\label{bl:cov:poles}
The tautological form and its residue satisfy
\begin{equation}\label{bl:cov:residue-form}
 \sigma\in H^0\bigl(Y,\omega_Y(R+D_h)\bigr),\qquad
 \Omega=\res_R(\sigma)\in H^0\bigl(R,\omega_R(D_h)\bigr).
\end{equation}
Here $\omega_R$ is the absolute dualizing sheaf of the reduced
simple-normal-crossing divisor $R$.  Both forms have weight $a$ and
deck character $\chi$.
\end{lemma}

\begin{proof}
Let $F$ be any prime divisor on $Y$, and restrict its valuation to a
function-field component of $Z$, obtaining $e v$.  On a downstairs
model carrying $v$, \eqref{bl:cov:tau} gives order
$q(A_{Z,\Delta_Z}(v)-1)$ for $\tau$.  Differential pullback adds
$q(e-1)$ to $e$ times this order.  Since $\sigma^{\otimes q}=\rho^*\tau$,
\begin{equation}\label{bl:cov:sigma-order}
 \ord_F(\sigma)+1=e A_{Z,\Delta_Z}(v).
\end{equation}
Log canonicity makes the right side nonnegative.  If the center lies
outside $E+H_Z$, it is strictly positive by the klt property.
The order on the left is integral, so $\sigma$ has at most a simple
pole, and every pole is over $E+H_Z$.  Because $E$ and $H_Z$ are
effective $\Q$-Cartier divisors, a prime whose center lies in this
union occurs in the support of their pullback.  That reduced support
is exactly $R+D_h$.  This proves the first assertion.  Adjunction for
the effective Cartier divisor $R$ on the smooth variety $Y$ gives
the residue sequence
\begin{equation}\label{bl:cov:residue-sequence}
 0\longrightarrow\omega_Y(D_h)
 \longrightarrow\omega_Y(R+D_h)
 \longrightarrow\omega_R(D_h)\longrightarrow0.
\end{equation}
It gives $\Omega$ and is equivariant, so residue preserves the stated
weight and character.
\end{proof}

The definition of $D_h$ is essential here.  An exceptional divisor
over $E\cap H_Z$ belongs to $D_h$ and is omitted from $R$, but its
possible simple pole is still allowed in~\eqref{bl:cov:residue-form}.
No disjointness condition on $H$ and the strata of $S$ has been used.

\subsection{The residue summand}

\begin{lemma}[Residue injection]\label{bl:lem:residue-injection}
Multiplication by $\Omega$ after pullback induces a split morphism
\begin{equation}\label{bl:cov:derived-split}
 \OO_E\longrightarrow \mathbf R(q_R)_*\omega_R(D_h)
\end{equation}
in the derived category of $\OO_E$-modules.  Consequently, for every
$i\geq0$, it induces an injection
\begin{equation}\label{bl:cov:injections}
 R^ig_*\OO_E\lhook\joinrel\longrightarrow
 R^ih_*\omega_R(D_h),\qquad e'\longmapsto\Omega e'.
\end{equation}
In particular this holds for $i=1$.  The map increases pullback
weight by $a$.
\end{lemma}

\begin{proof}
\emph{The character summand.}
We first identify the $\chi$-eigensheaf:
\begin{equation}\label{bl:cov:eigensheaf}
 \bigl(\rho_*\omega_Y(D_h)\bigr)_\chi
   =\sigma\cdot\OO_Z(-E).
\end{equation}
A meromorphic $\chi$-form divided by $\sigma$ is deck-invariant, so it
is the pullback of a meromorphic function $v$ on $Z$, by the invariant
algebra calculation following~\eqref{bl:cov:second-algebra}.  At a prime
of $Z$ with coefficient $d$ in $\Delta_Z$, and at a prime of the cover
above its generic point, \eqref{bl:cov:sigma-order} becomes
\[
 \ord_F(v\sigma)=e\ord(v)+e(1-d)-1.
\]
Over $E$ no pole is permitted in $\omega_Y(D_h)$, so $d=1$ forces
$\ord(v)\geq1$.  Over $H_Z$ a simple pole is permitted and forces
$\ord(v)\geq0$.  Elsewhere $0\leq d<1$, and the integer
$e(1-d)$ lies between $1$ and $e$; hence holomorphy is equivalent
to $\ord(v)\geq0$.  These necessary inequalities, and normality,
give $v\in\OO_Z(-E)$.  Conversely such a function pulls back to a
regular function vanishing along $\widehat E$.  Since
$\widehat E\geq R$, it cancels every pole of $\sigma$ outside $D_h$
on the whole resolution, including exceptional divisors.  This
proves~\eqref{bl:cov:eigensheaf}.

\smallskip\noindent
\emph{Vanishing and the thickened residue.}
We next prove the required vanishing:
\begin{equation}\label{bl:cov:vanishing}
 R^j\rho_*\omega_Y(D_h)=0\qquad(j>0).
\end{equation}
Choose a rational $\epsilon>0$ so small that every nonzero coefficient
of $F=\epsilon\rho^*H_Z$ lies strictly between $0$ and $1$.
Then $D_h-F$ is an effective simple-normal-crossing boundary with
all coefficients less than one, and
\[
 (K_Y+D_h)-\bigl(K_Y+(D_h-F)\bigr)=F.
\]
The divisor $F$ is $\rho$-nef, being a pullback from $Z$, and
$\rho$-big: $\rho$ is projective and generically finite on every
component, and every divisor restricts to a big divisor on its
zero-dimensional generic fiber.  Relative Kawamata--Viehweg vanishing
therefore applies to the Cartier divisor $K_Y+D_h$ and the klt
boundary $D_h-F$; see \cite[Theorem~3.3]{FujinoVanishing}.
If $H=0$, take $F=D_h=0$; the same relative bigness observation
includes this case.  This proves~\eqref{bl:cov:vanishing} without
asserting any absolute nefness of an auxiliary adjoint.

Since $E$ is Cartier, the projection formula gives the same vanishing
for $\omega_Y(\widehat E+D_h)$ and, by
\eqref{bl:cov:eigensheaf}, identifies their derived eigensummands as
\begin{equation}\label{bl:cov:two-derived-images}
 \begin{split}
 \bigl(\mathbf R\rho_*\omega_Y(D_h)\bigr)_\chi
     &\simeq\OO_Z(-E),\\
 \bigl(\mathbf R\rho_*\omega_Y(\widehat E+D_h)\bigr)_\chi
     &\simeq\OO_Z.
 \end{split}
\end{equation}
Both are concentrated in degree zero; the isomorphisms are
multiplication by $\sigma$.  The inclusion of the two sheaves on $Y$
corresponds to the ordinary ideal inclusion $\OO_Z(-E)\hookrightarrow
\OO_Z$.

For the possibly nonreduced effective Cartier divisor $\widehat E$,
adjunction still gives
\[
 \omega_{\widehat E}(D_h)
   =\omega_Y(\widehat E+D_h)/\omega_Y(D_h).
\]
Taking the quotient in~\eqref{bl:cov:two-derived-images} shows that
\begin{equation}\label{bl:cov:thick-summand}
 \bigl(\mathbf R(q_{\widehat E})_*
                 \omega_{\widehat E}(D_h)\bigr)_\chi
       \simeq\OO_E
\end{equation}
in the derived category on $E$.  Here this conclusion can be checked
on cohomology sheaves after the exact faithful pushforward by the
closed immersion $E\hookrightarrow Z$: the higher eigencohomology
vanishes and the degree-zero map is the residue of multiplication by
$\sigma$.  There is no need to assert full faithfulness of derived
pushforward for that immersion.

\smallskip\noindent
\emph{Retraction from the reduced residue.}
We construct the splitting map on $E$ itself.
Pullback of functions and multiplication by $\Omega$ give
\[
 \alpha:\OO_E\longrightarrow\mathbf R(q_R)_*\omega_R(D_h).
\]
The inequality $R\leq\widehat E$ gives an inclusion of residue
quotients on $Y$, hence a morphism
\[
 \beta:\mathbf R(q_R)_*\omega_R(D_h)
       \longrightarrow\mathbf R(q_{\widehat E})_*
                                      \omega_{\widehat E}(D_h).
\]
These are $\OO_E$-linear: multiplication by the pulled-back equation
$y$ kills both quotients.  If a local function on $E$ is lifted to
$v\in\OO_Z$, the composite $\beta\alpha$ is represented by the
class of $\rho^*v\,\sigma$ modulo $\omega_Y(D_h)$.
Changing $v$ by a multiple of $y$ does not change this class, since
$y\sigma\in\omega_Y(D_h)$ by $\widehat E\geq R$.
Thus projection onto the $\chi$-summand followed by
\eqref{bl:cov:thick-summand} sends $\beta\alpha$ to the identity.
The projection is available in the derived category because
averaging over $\mu_q$ gives the idempotent
$q^{-1}\sum_{\zeta\in\mu_q}\chi(\zeta)^{-1}\zeta^*$.
This supplies a retraction of $\alpha$ and proves
\eqref{bl:cov:derived-split}.  Finally apply $\mathbf Rg_*$ and take
cohomology to obtain~\eqref{bl:cov:injections}.  The weight statement
follows from $\wt(\Omega)=a$.
\end{proof}

\subsection{The obstruction to be detected}
\label{bl:cov:obstruction-contract}

Put $A_d=|g|_*(\OO_Z/I^d)$ for $d\geq1$, pushing the underlying
sheaves of complex vector spaces in the Zariski topology along
$|g|:|E|\to|B|$. This uses no morphism from a thickening to $B$.
Fix $k\geq1$ and suppose that $A_{d+1}\to A_d$ is surjective
for $1\leq d<k$.
The current connecting map is
$\partial_k:A_k\to R^1g_*(I^k/I^{k+1})$.
Section~\ref{bl:sec:lifting} will show that it factors, after framing
$I^k/I^{k+1}$ by $y^k$, through a derivation
\[
 \delta_k:g_*\OO_E\longrightarrow R^1g_*\OO_E,
 \qquad \partial_k(\widetilde F)=y^k\delta_k(F).
\]
Here $\widetilde F\in A_k$ is any local lift of $F\in g_*\OO_E$.
The factorization and derivation rule will be proved there from the
previous surjections. Since $y$ has weight $-1$, this derivation
has degree $k$ for the frame action.

Restrict $\delta_k$ to pullbacks of functions on $B$. In local
coordinates $t_i$, the residue injection produces a global section
of $R^1h_*\omega_R(D_h)\otimes T_B$, given by
\[
 v_k=\sum_i\Omega\delta_k(g^*t_i)\otimes\partial_{t_i}
 ,\qquad \wt(v_k)=a+k>0.
\]
The derivation rule makes this tensor independent of coordinates.
The next two sections develop a differential-operator test for it.
The graph identity that makes the test applicable, and that then
annihilates $\delta_k$ on all of $g_*\OO_E$, is proved with the
induction in Section~\ref{bl:sec:lifting}.
